% !TEX root = ../main.tex
% =============================================================================
% APPENDIX - implementation settings needed to reproduce the downstream
% experiments.
%
% The Appendix heading and numbering are defined in main.tex.
% This file starts at the section level.
%
% VOICE: impersonal and concise.
% =============================================================================

This appendix provides the implementation settings used for the downstream
experiments reported in Chapter~\ref{ch:results}.

\section{Downstream Adaptation}\label{app:adaptation}

Table~\ref{tab:app_adaptation} shows the optimisation settings used for
linear probing and finetuning. Within each adaptation setting, the same
hyperparameters were used for all four downstream tasks, all three
BrainLM-EEG variants, and all evaluation protocols.

\begin{table}[htbp]
    \centering
    \caption[Downstream adaptation settings]{Optimisation settings for linear
    probing and finetuning.}
    \label{tab:app_adaptation}

    \begin{tabularx}{\textwidth}{
        >{\raggedright\arraybackslash}p{4.6cm}
        >{\raggedright\arraybackslash}X
        >{\raggedright\arraybackslash}X}
        \toprule
        \textbf{Setting} &
        \textbf{Linear probing} &
        \textbf{Finetuning} \\
        \midrule

        Encoder parameters & Frozen & Updated \\
        Learning rate & \(1\times10^{-3}\) & \(5\times10^{-5}\) \\
        Weight decay & \(0\) & \(0.05\) \\
        Learning rate schedule & Cosine decay &
        Linear warmup over 10 epochs, followed by cosine decay \\
        Maximum epochs & 50 & 50 \\
        Early stopping patience & 10 & 10 \\
        Batch size & 128 & 128 \\
        Gradient clipping (maximum norm) & \(1.0\) & \(1.0\) \\

        \bottomrule
    \end{tabularx}
    \\[3pt]
    \begin{minipage}{\linewidth}\footnotesize\raggedright
        All runs used the AdamW optimiser
        \parencite{loshchilovDecoupledWeightDecay2019}, a random seed of 42,
        and cross-entropy loss. Early stopping and checkpoint selection were
        based on validation loss.
    \end{minipage}
\end{table}

The classification head described in Section~\ref{sec:m_lp_ft} was a
three-layer perceptron:

\begin{center}
CLS \(\rightarrow\) Linear(512, 1024) \(\rightarrow\) GELU \(\rightarrow\)
Dropout(0.1) \(\rightarrow\)\\
Linear(1024, 1024) \(\rightarrow\) GELU \(\rightarrow\) Dropout(0.1)
\(\rightarrow\) Linear(1024, 2)
\end{center}

The same classification head was used in both settings.
The encoder remained frozen during linear probing and was updated during finetuning.

All downstream classification results were obtained from the CLS token.
The layer-wise analysis in Section~\ref{sec:m_layer_probe} was the only
exception. For this analysis, the patch token representations from each
encoder stage were averaged, and the CLS token was excluded. Therefore,
the results in Section~\ref{sec:r_layers} are not directly comparable with
the CLS-based results reported elsewhere.

\section{Electrode Montage and Coordinates}\label{app:montage}

Channel names were standardised to the current electrode nomenclature.
The legacy labels T3, T4, T5, and T6 were mapped to T7, T8, P7, and P8,
respectively.

The channels were ordered from anterior to posterior and from left to right
within each row. This order was kept fixed so that each spatiotemporal token
corresponded to the correct electrode.

The azimuth and elevation angles used by BrainLM-EEG-RoPE
(Section~\ref{sec:m_rope}) were derived from the standard 10-20 electrode
coordinates provided by \texttt{MNE-Python}. The coordinates use a
head-centred Cartesian system. The \(x\) axis points towards the right ear,
the \(y\) axis points towards the nasion, and the \(z\) axis points upwards.

Each electrode position was projected onto the unit sphere. Azimuth and
elevation were then calculated as

\begin{equation}
    \theta_c = \operatorname{atan2}(x_c, y_c),
    \qquad
    \varphi_c = \arcsin(z_c)
    \label{eq:app_rope_angles}
\end{equation}

Table~\ref{tab:app_montage} lists the electrode order and the corresponding
coordinates.

\begin{table}[htbp]
    \centering
    \small
    \setlength{\tabcolsep}{4pt}
    \caption[Common electrode montage and coordinates]{The 19 common electrodes
    in the order used by the model, with their spherical coordinates.}
    \label{tab:app_montage}

    \begin{tabular}{c l c c @{\hspace{2em}} c l c c}
        \toprule
        \multirow{2}{*}{\textbf{Index}} &
        \multirow{2}{*}{\textbf{Electrode}} &
        \textbf{Azimuth} &
        \textbf{Elevation} &
        \multirow{2}{*}{\textbf{Index}} &
        \multirow{2}{*}{\textbf{Electrode}} &
        \textbf{Azimuth} &
        \textbf{Elevation} \\

        & & \(\theta_c\) & \(\varphi_c\) &
        & & \(\theta_c\) & \(\varphi_c\) \\
        \midrule

         0 & Fp1 & \(-0.337\) & \(-0.078\) & 10 & C4 & \(+1.732\) & \(+0.752\) \\
         1 & Fp2 & \(+0.338\) & \(-0.079\) & 11 & T8 & \(+1.746\) & \(-0.109\) \\
         2 & F7  & \(-1.027\) & \(-0.138\) & 12 & P7 & \(-2.363\) & \(-0.024\) \\
         3 & F3  & \(-0.758\) & \(+0.523\) & 13 & P3 & \(-2.549\) & \(+0.532\) \\
         4 & Fz  & \(+0.005\) & \(+0.849\) & 14 & Pz & \(+3.138\) & \(+0.795\) \\
         5 & F4  & \(+0.762\) & \(+0.498\) & 15 & P4 & \(+2.525\) & \(+0.531\) \\
         6 & F8  & \(+1.024\) & \(-0.139\) & 16 & P8 & \(+2.356\) & \(-0.025\) \\
         7 & T7  & \(-1.759\) & \(-0.109\) & 17 & O1 & \(-2.886\) & \(+0.076\) \\
         8 & C3  & \(-1.747\) & \(+0.770\) & 18 & O2 & \(+2.882\) & \(+0.076\) \\
         9 & Cz  & \(+3.098\) & \(+1.480\) &    &    &            &            \\

        \bottomrule
    \end{tabular}
    % No blank line before \\[3pt] -- see the note on the table above.
    \\[3pt]
    \begin{minipage}{\linewidth}\footnotesize\raggedright
        Azimuth \(\theta_c\) and elevation \(\varphi_c\) are given in radians.
        These values were used directly as the scalar positions in
        Equation~\ref{eq:rope1d}. Rotation frequencies used the base
        \(b=10^{4}\), without additional rescaling.
    \end{minipage}
\end{table}

\FloatBarrier

\section{Code Availability}\label{app:code}

The implementation, configuration files, and scripts used in this thesis are
publicly available at \url{https://github.com/gokul-pv/ssl-eeg}.

% TODO: cite a tagged release or commit hash for the final thesis version.
